\documentclass[10pt,twocolumn,letterpaper]{article}
\usepackage[rebuttal]{cvpr}

% Include other packages here, before hyperref.
\usepackage{graphicx}
\usepackage{amsmath}
\usepackage{amssymb}
\usepackage{booktabs}
\usepackage{url}
\usepackage{microtype}

\usepackage[pagebackref,breaklinks,colorlinks,bookmarks=false]{hyperref}

% Support for easy cross-referencing
\usepackage[capitalize]{cleveref}
\crefname{section}{Sec.}{Secs.}
\Crefname{section}{Section}{Sections}
\Crefname{table}{Table}{Tables}
\crefname{table}{Tab.}{Tabs.}

%%%%%%%%% PAPER ID  - PLEASE UPDATE
\def\cvprPaperID{} 
\def\confName{MAI}
\def\confYear{Fall 2026}

\begin{document}

%%%%%%%%% TITLE - PLEASE UPDATE
\title{Paper title \\ Your name and student ID}  % **** Enter the paper title here

\maketitle
\thispagestyle{empty}
\appendix

%%%%%%%%% BODY TEXT - ENTER YOUR RESPONSE BELOW
\section{Summary}
Summarize in your own words the paper, its key ideas, its contributions and their significance. This summary helps the AC and the authors understand the rest of your review and be confident that you understand the paper.

\section{Strengths}
\begin{enumerate}
    \item Detail the strengths of the paper.
    \item These can be for example: interesting ideas, an insightful organization of related work, new tools, impressive results.
    \item Be specific: Your comments will be much more helpful to the ACs and the authors than your score.
\end{enumerate}

\section{Major Weaknesses}
\begin{enumerate}
    \item Detail the important weaknesses you identified and justify them. For example:
    \item If you found the method difficult to understand, explain what the authors did wrong and give examples.
    \item If the claimed contributions are not backed up by the experiments, explain what should have been done.
    \item Be specific: Do not simply give summary judgments. Writing “not novel”, “unclear”, “incorrect” is not sufficient!
    \item Your comments will be much more helpful to the ACs and the authors than your score.
    \item (2026) Suppose this paper was submitted in 2026, what would be its major weaknesses from today’s point of view?
    \item (2026) In other words, this comment is based on current technological progress, not the time when the paper was originally published.
\end{enumerate}

\section{Minor Weaknesses}
\begin{enumerate}
    \item Mention the weaknesses that are easy to fix in a revision of the paper. For example:
    \item Missing references that should be discussed but do not really change the novelty of the paper.
    \item Occasional typographic errors.
    \item Minor unclear points.
    \item Mention here any minor suggestions, questions, corrections, etc. that can help the authors improve their paper.
    \item (2026) Suppose this paper was submitted in 2026, what would be its minor weaknesses from today’s point of view?
    \item (2026) In other words, this comment is based on current technological progress, not the time when the paper was originally published.
\end{enumerate}

\section{Justification}
Balance the strengths and the weaknesses to explain why the paper should be accepted or not. Does the paper present sufficient knowledge advancement?  Use your best judgment.

(2026) Suppose this paper was submitted in 2026, on what grounds would you recommend accepting or rejecting it from today’s point of view? In other words, this comment is based on current technological progress, not the time when the paper was originally published.

\section{Final Rating}
5: Strong Accept

4: Weak Accept

3: Borderline

2: Weak Reject

1: Strong Reject

Provide (2026) final rating as well.

This template follows the review format of CVPR 2025 \cite{web} and \cite{slides}.

\section{Reflections}
Share your reflections on this paper. What did you learn (technologies, writing, etc.)? How did it inspire you? Is it useful for your project or research? How might you apply it to your research? Just examples. Feel free to connect the paper to yourself in other ways.



%%%%%%%%% REFERENCES
{\small
\bibliographystyle{ieee_fullname}
\bibliography{egbib}
}

\end{document}
